% 章节依据：陈维桓《微分几何》第2版 LaTeX 原文。
% 仿照数理统计“期末笔记.tex”：按能力与功能组织，再列所用工具。
% 能力分组是笔记的学习路线，并非教材原章名；同一章可服务多项任务。
\chapter*{微分几何能解决什么问题}
\addcontentsline{toc}{chapter}{微分几何能解决什么问题}

\begingroup
\setstretch{1.10}
\dirtree{%
.1 \textsf{微分几何的功能与作用}\quad 问题 $\longrightarrow$ 工具 $\longrightarrow$ 结果.
.2 \textsf{基础工具层：把几何对象变成可计算的数据}.
.3 描述位置与变化 $\;\to\;$ 坐标、标架、向量函数与导数（第1章）.
.3 描述曲线、曲面 $\;\to\;$ 参数方程、切向量、切平面与法向量（第2、3章）.
.2 \textsf{能力一：测量长度、角度和面积}.
.3 曲线有多长 $\;\to\;$ 弧长积分与弧长参数（第2章）.
.3 曲面上怎样测量 $\;\to\;$ 第一基本形式 $\mathrm{I}$，计算长度、夹角与面积（第3章）.
.2 \textsf{能力二：判断对象怎样弯曲}.
.3 曲线怎样弯、怎样扭 $\;\to\;$ 曲率 $\kappa$、挠率 $\tau$ 与 Frenet 公式（第2章）.
.3 曲面沿不同方向怎样弯 $\;\to\;$ 第二基本形式 $\mathrm{II}$、法曲率与主曲率（第4章）.
.3 曲面内部的弯曲怎样测 $\;\to\;$ 由 $\mathrm{I}$ 确定高斯曲率 $K$（第5、6章）.
.2 \textsf{能力三：比较形状，并从几何数据重建对象}.
.3 比较曲面上的测量规则 $\;\to\;$ 保长对应、保角对应（第3章）.
.3 判断曲线是否全等、构造曲线 $\;\to\;$ 曲率、挠率与曲线论基本定理（第2章）.
.3 判断曲面是否全等 $\;\to\;$ 第一、第二基本形式与曲面唯一性定理（第5章）.
.3 给定数据能否成为曲面 $\;\to\;$ 检查 Gauss--Codazzi 相容条件，再局部重建（第5章）.
.2 \textsf{能力四：在曲面上寻找路径，并沿路径比较方向}.
.3 路径在曲面内怎样转弯 $\;\to\;$ 测地曲率；为零时为测地线（第6章）.
.3 怎样求最短路径的候选 $\;\to\;$ 测地线方程，再判断最短性条件（第6章）.
.3 不同点的切向量怎样比较 $\;\to\;$ 协变导数与沿路径的平行移动（第6章）.
.2 \textsf{能力五：用局部曲率研究整体性质}.
.3 求测地三角形的角和 $\;\to\;$ Gauss--Bonnet 公式（第6章）.
.3 联系总曲率与拓扑 $\;\to\;$ 紧致无边可定向曲面上 $\int_S K\,\dd A=2\pi\chi(S)$（第6章）.
.2 \textsf{工具升级：把上述计算写成统一的方程}.
.3 跟踪切向与法向的变化 $\;\to\;$ 活动标架、微分形式与结构方程（第7章）.
.3 处理曲面与曲面上的曲线 $\;\to\;$ 用外微分法组织曲率和相容性计算（第7章）.
}
\endgroup

\vspace{0.8em}
\noindent\textsf{章节如何配合：}
第1章提供计算工具；第2章同时服务测量、弯曲与重建；
第3章提供曲面测量规则，第4章补充空间中的弯曲数据，第5章据此比较和重建曲面；
第6章用这些量研究路径及整体性质；第7章提供统一计算方法。
